\section*{Abstract}
The magnitude/phase decomposition of the FFGFT mass operator (Doc.~288/289) reorganises the information
question. Information has a \emph{magnitude side} (mass/energy -- Landauer and Vopson, derived by FFGFT
from the $T^4$ geometry) and an independent \emph{phase side} (holonomy, mixing -- the one open datum
$\theta$). The cleanest witness is the \textbf{photon}: exactly massless -- zero rest-mass magnitude, hence pure phase in the mass channel -- and the quantum of
the $U(1)$ gauge potential $A_\mu$, which is the connection; FFGFT's information unit $E_\text{bit}=\hbar c/L$ \emph{is} a photon
energy. The \emph{amount} of information is not representation-relative but absolute and dimensionless,
anchored in $\xi$: $I_1=\log_2(1/\xi)\approx12.87$ bit per degree of freedom; energy and capacity are its
dimensional, frequency-scaling dressings. The question is not \enquote{is information fundamental?} but
\textbf{which channel} -- and only the geometry holds both. A structural diagnosis, not a new derivation.

\section{Starting point: geometry first, and the magnitude/phase lens}
FFGFT's stance on the information question is documented and geometry-first. Doc.~243 (discussion with
Onur Teker, UIFT): the map of FFGFT into an information formalism is a \emph{projection} $\pi$, not an
isomorphism; the bit count $I=\log_2|\mathcal{W}|$ \emph{presupposes} the geometric state space.
Doc.~255/257 derive Landauer's principle and Vopson's mass-information equivalence \emph{from} the $T^4$
geometry: a bit is a minimal topological winding $\Delta w=1$, \enquote{information is topologically
frozen kinetic energy}, $E_\text{bit}(L)=\hbar c/L=pc$. Doc.~274 places FFGFT \emph{geometry-primitive}
beside information-primitive (Vopson). FFGFT does not reject Vopson -- it grounds him.

The new lens is the magnitude/phase decomposition of the mass operator (Doc.~288/289): the operator is a
Z$_3$-circulant, DFT-diagonalised, and its content splits into a \emph{magnitude spectrum} (masses,
energy) and a \emph{phase spectrum} (mixing, holonomy). The $\theta$ finding (Doc.~282) says the
magnitude spectrum fixes the phase spectrum only up to an all-pass factor (made precise in §2). That is the key to all that follows.

\section{Two channels: magnitude and phase}
$\Delta w=1$ and $E_\text{bit}=\hbar c/L$ are magnitude/energy statements. Since the magnitude spectrum
fixes the phase spectrum only up to an all-pass factor (see below):

Vopson's mass-information equivalence -- which FFGFT derives -- structurally captures \emph{only the
magnitude channel}. The phase channel (holonomy, mixing, $\theta$) is independent information, not
reducible to mass/energy.

This shows concretely in two places. \emph{First, mixing-blindness:} a mass operator
$M=U\,\mathrm{diag}(m)\,U^{\dagger}$ carries the eigenvalues $m$ (magnitude channel) and the eigenvectors
$U$ (mixing, phase channel); any function of the eigenvalues alone ($\sum m$, Koide $Q$) is invariant
under $U$. Numerically (\texttt{masse\_information\_phasenblind.py}): for fixed masses $(0.5,1.5,3.0)$,
$\sum m=5$ and $Q=0.372$ stay equal across three mixings while $\lVert M_\text{offdiag}\rVert$ ($1.42$ to
$1.78$) varies -- two physically distinct operators with the same mass-information.

\emph{Second, the phase carries the structure:} known in signal engineering since Oppenheim/Lim (1981) --
swapping the magnitude and phase spectra of two patterns, the reconstruction follows the \emph{phase}
donor. Numerically (\texttt{phase\_traegt\_struktur.py}): phase-only correlates with the original
($+0.53$), magnitude-only does not ($-0.24$). A mass-energy accounting thus captures the
\emph{structurally poorer} channel.

\emph{Third, precisely:} \enquote{the magnitude does not fix the phase} holds only \emph{up to an
all-pass factor}. Every stable transfer function factors as $H=H_\text{min}\cdot H_\text{ap}$: for the
\emph{minimum-phase} part the magnitude does fix the phase (Bode/Hilbert relation,
$\arg=\mathcal{H}\{\log|H|\}$), whereas the \emph{all-pass} ($|H_\text{ap}|=1$) carries a phase the
magnitude does \emph{not} fix. Numerically (\texttt{minimalphase\_allpass.py}): the phase reconstructed
from the magnitude alone matches $\arg H_\text{min}$ (error $\sim10^{-15}$), the residual is exactly the
all-pass. Free -- and thus the open part -- is the \emph{all-pass}, and that is where the $\theta$
holonomy sits (the minimum-phase\,(FFGFT)\,$\times$\,all-pass\,(HLV-$R$) reading, Doc.~288).

\section{The photon: the exact anchor}
The cleanest witness is the photon -- and the corpus already holds the pieces. Doc.~007 treats neutrinos
via the photon analogy: photon $E^2=(pc)^2+0$ (exactly massless), neutrino
$E^2=(pc)^2+(\sqrt{\xi^2/2}\,mc^2)^2$ (quasi-massless). The photon is the $m=0$ limit, the neutrino
\enquote{photon plus the first $\xi^2/2$ correction}.

The photon has rest-mass magnitude exactly zero -- the exact all-pass \emph{in the mass channel}: all of its rest-mass content is phase. (Its energy $\hbar\omega$ remains a magnitude quantity; \enquote{pure phase} refers to the rest-mass, not the energy channel.) Three
things coincide:
\begin{enumerate}
  \item \textbf{Unfrozen information.} The photon is the unfrozen case of \enquote{information $=$ frozen
        kinetic energy} ($E=pc$); FFGFT's unit $E_\text{bit}=\hbar c/L$ \emph{is} a photon energy,
        Landauer's thermal bit its freezing.
  \item \textbf{Photon and connection.} The gauge potential $A_\mu$ \emph{is} the $U(1)$ connection; the
        photon is its quantum. \emph{One} gauge-invariant observable is the holonomy ($\oint A\,\mathrm{d}l$,
        Aharonov-Bohm, even where $F=0$) -- alongside energy $\hbar\omega$, momentum and polarization,
        which are not holonomies. The link to \enquote{it from connection} is \emph{conceptual}
        (relational character), not a formal identity with $\theta$.
  \item \textbf{Counterexample to \enquote{mass $=$ information}.} Zero rest mass, yet carrier of
        essentially all our information transmission.
\end{enumerate}

The neutrino is only the first correction to this: a \enquote{damped photon} with a tiny magnitude
instead of exactly zero. In the degenerate limit $M=m\cdot\mathbb{1}$ for any $U$
($\lVert M-m\mathbb{1}\rVert\sim10^{-16}$), the mixing would leave $M$, the oscillation content sitting in
the propagation/holonomy phase -- the same situation the photon realises \emph{exactly} and without a
degeneracy hypothesis.

\section{What \enquote{amount of information} means: packet, count, energy, capacity}
Is a photon the smallest packet? Only \emph{per mode}. A single-frequency photon is a plane wave; a
localised photon is a wave packet. Crucially: a time-limited packet keeps the
\emph{definite} carrier $f_0$ (the laser colour, the energy jump); the limitation only produces a band
$\Delta f\sim1/\Delta t$ \emph{around} $f_0$ (window effect), not a spectrum. The cycle count is
$N=f_0\,\Delta t$: at fixed time window $\propto f_0$, but at fixed cycle count the fractional bandwidth
$\Delta f/f_0=1/N$ is \emph{frequency-independent}. For a photon from an energy jump, the lifetime $\tau$
is the window: the natural linewidth $\Delta f=f_0/Q$ with $Q=\omega_0\tau=2\pi N$ ($\tau=10$\,ns,
$f_0=5\cdot10^{14}$\,Hz $\Rightarrow N=5\cdot10^6$, $\Delta f\approx16$\,MHz; numerically
\texttt{fenster\_zyklen\_linienbreite.py}). Definite carrier \emph{plus} window gives the band.

\paragraph{Neither cycle count nor window length is intrinsic.} The emission process sets both, not
\enquote{photon-hood}: a monochromatic mode has infinitely many cycles and an infinite window (plane
wave); a spontaneously emitted photon inherits $f_0$ \emph{and} $\tau$ from the transition
($N=f_0\tau=Q/2\pi$, the two coupled); a technical packet fixes one quantity, the other follows from
$N=f_0\Delta t$. Intrinsic -- scale-invariant -- is the dimensionless number $N$ ($\Delta f/f=1/N$);
$\Delta t=N/f$ is its dimensional dressing. The span is large: few-cycle pulse $N\approx2$, narrow atomic
line $N\approx5\cdot10^6$ (numerically \texttt{quant\_bit\_modeninfo.py}) -- there is no universal cycle
count and no universal window length.

Three quantities then separate cleanly (numerically \texttt{photon\_paket\_count\_energie.py}):

\begin{itemize}
  \item \textbf{Count}: the quantum number (occupation) is a \emph{count} -- one excitation,
        frequency-independent, relational --, \emph{not} a bit measure. What a quantum carries as
        \emph{information} is below.
  \item \textbf{Energy} per bit: $E=\hbar\omega=\hbar c/L$ -- frequency-dependent (radio
        $\sim2\cdot10^{-10}$\,eV to gamma $\sim2\cdot10^{5}$\,eV).
  \item \textbf{Capacity}: Shannon $C=B\log_2(1+\mathrm{SNR})$, $B\propto f$ -- frequency-dependent
        ($\sim2\cdot10^{5}$ to $\sim2\cdot10^{20}$ bit/s).
\end{itemize}

\enquote{Higher frequency $\to$ more information} is exactly the capacity: more bandwidth, more modes per
second. That is a \emph{different} axis from magnitude/phase: \emph{within} a mode each symbol carries
amplitude \emph{and} phase (I/Q, QAM), \emph{across} modes the bandwidth counts. FFGFT mirrors both:
within the mode the Z$_3$-circulant ($r,\theta$), across modes $|\mathcal{W}|$.

\paragraph{A quantum is a count, not a bit.} \enquote{One quantum $=$ one bit} is not a general equation
-- it confuses counting with measuring. A quantum is one \emph{excitation} (a count of~$1$); its
\emph{information content} is $\log_2 M$ bit when it is chosen from $M$ distinguishable modes --
\emph{variable}, capped per degree of freedom at $I_1=\log_2(1/\xi)\approx12.9$ bit (§5). \enquote{One
quantum $=$ one bit} holds only for $M=2$ (the mode fixed to a binary alphabet) -- a coding convention,
not a law. Count and bit content are different axes: one quantum among $2^{13}$ modes carries $13$~bit,
and so do $13$ binary quanta -- same bit count, different quantum number (numerically
\texttt{quant\_bit\_modeninfo.py}). The quantum \emph{counts}; $\log_2 M\le I_1$ \emph{measures}. This
also places §5 cleanly: $I_1$ is not the content of \emph{one} quantum but the absolute \emph{ceiling} per
degree of freedom, which the variable mode content $\log_2 M$ never exceeds.

\section{The absolute amount: anchored in $\xi$}
An earlier reading -- information content is representation-relative, there is no absolute bit count --
was overstated and unfaithful to FFGFT. The absolute, frequency-independent content is the
\emph{dimensionless} quantity, and $\xi$ anchors it. $\xi=4/30000$ is dimensionless (the absolute
parameter); $\hbar,c$ are SI bookkeeping.

$\xi$ caps the winding number at $n_\text{max}\sim1/\xi$ (Doc.~243/009) -- without this bound
$I=\log_2|\mathcal{W}|$ diverges. One winding coordinate can take $\sim1/\xi$ distinguishable values:
$I_1=\log_2(1/\xi)=\log_2(7500)\approx12.87$~bit. Total
$I_\text{total}\sim N_\text{max}\cdot\log_2(1/\xi)$ (cells $N_\text{max}=V/L_0^{d}$, Doc.~251, times
depth). Dimensionless, set by $\xi$ -- the energy $\hbar c/L$ is not the amount but its dimensional price,
and the capacity saturates at the $\xi$ scale.

Only \emph{weakly} relative are the bit count of a \emph{description} at a chosen resolution and what a
lossy projection discards (Doc.~243) -- a statement about approximations, not about the information
itself. The full information (the geometry, or its lossless Hilbert image, Doc.~230) is absolute.

\paragraph{Absolute and relational are not in conflict.} The two statements -- the amount is
\emph{absolute} (§5), the open datum $\theta$ is \emph{relational} (§8) -- lie on different axes:
\enquote{absolute} measures the \emph{amount}, \enquote{relational} describes the \emph{nature} of a
datum. A holonomy (winding) is the paradigm of \emph{both}: relationally defined (a loop integral, not a
local point value) \emph{and} absolutely valued (gauge-invariant, integer). Numerically
(\texttt{absolut\_relational.py}): a winding $w=3$ is readable only as a loop integral (the half loop
gives $1$) and stays unchanged under a single-valued rephasing (gauge-invariant); and the absolute number
$\log_2(1/\xi)$ \emph{counts} exactly such relational objects ($\sim15000$ winding sectors up to
$n_\text{max}\sim1/\xi$, $\log_2\approx13.9$ bit). The absoluteness sits in the ground (the geometry,
$\xi$), the relationality in the degrees of freedom upon it -- the absolute bit count counts relational
windings.

\section{Static and dynamic definition of the bit}
The definition so far is throughout \emph{static}: the bit as structure -- winding $\Delta w=1$, counting
ceiling $\log_2 M\le I_1$, energy $E_0=\hbar c/L$ (magnitude side). That says \emph{what} a bit is, not
what it \emph{does}. The dynamic definition is not an add-on but the conjugate side -- and it itself splits
into two, in which energy plays \emph{different} roles.

\emph{Kinematic (reversible), the $\tilde T$ side:} a flip takes $\tilde T=\hbar/E_0=1/m_0$, the rate is
$1/\tilde T=E_0/\hbar$. Here $E_0\cdot\tilde T=\hbar$ is nothing but $\tilde T\cdot m=1$ -- the time-mass
duality, the FFGFT foundation. The quantum speed limit (Margolus-Levitin) $\tau\ge\pi\hbar/2E$ sharpens it:
energy \emph{sets the clock} without being dissipated (unitary). The cost is \emph{time}. The capacity axis
(bit/s, §4) is exactly this rate.

\emph{Thermodynamic (irreversible), the Landauer side:} a bit carries entropy $k_B\ln 2$; \emph{erasing}
it (forgetting) into a bath at temperature $T$ dissipates $k_B T\ln 2$ of heat. This depends on $T$,
\emph{not} on $E_0$, and the heat is \emph{dissipated} -- the cost is \emph{heat}, not time. In FFGFT this
side is \emph{emergent}: the fundamental $T^4$ dynamics is unitary-reversible, so the kinematic
($\tilde T$) face is primary and the Landauer price arises only through coarse-graining, a bath, and
logical irreversibility.

The dimensionless count $\log_2 M$ is primary; the constants dress it (numerically
\texttt{bit\_statisch\_dynamisch.py}): $\hbar,c$ into energy $E_0$ and time $\tilde T$ (kinematic,
reversible: \emph{time} cost, $\tilde T\cdot m=1$), $k_B$ into entropy $k_B\ln 2$ and erase heat
$k_B T\ln 2$ (thermodynamic, irreversible: \emph{heat} cost). $\hbar,c,k_B$ are SI bookkeeping. Energy
plays two roles: \emph{clock} (kinematic, no loss) versus \emph{dissipated heat} (thermodynamic).

Important: \enquote{bookkeeping} does not mean \emph{dispensable}. $k_B,\hbar,c$ are the \emph{necessary}
unit bridges. SI separates temperature (K), energy (J), time (s) into their own units, and as soon as one
\emph{computes} thermodynamically -- and that happens in SI -- $k_B$ is the indispensable conversion
K$\to$J (otherwise the erase heat would stand in kelvin rather than joule), just as $\hbar,c$ carry the
conversion to $E_0=\hbar c/L$. In natural units ($k_B=\hbar=c=1$) the costs are pure numbers ($\ln 2$); SI
introduces the unit walls, and the bookkeeping is their mortar -- ontologically secondary, operationally
indispensable (numerically \texttt{buchhaltung\_si\_natuerlich.py}).

\paragraph{The bit in SI numbers.} Anchored at the framework's characteristic energy
$E_0=\sqrt{m_e m_\mu}=7.348$\,MeV (coupled to the fine-structure constant via $\alpha=\xi E_0^2$; the
variant $\sqrt{\alpha/\xi}=7.398$\,MeV differs by $+0.68\%$ -- not exact equality). The generic
$E_0=\hbar c/L$ coincides with this $E_0$ exactly at $L=L_0=\hbar c/E_0$. Then (numerically
\texttt{bit\_in\_si.py}): $E_0=1.177\cdot10^{-12}$\,J, $L_0=2.686\cdot10^{-14}$\,m $=26.85$\,fm,
$m_0=1.310\cdot10^{-29}$\,kg, $\tilde T=\hbar/E_0=8.958\cdot10^{-23}$\,s $=L_0/c$ and
$1/\tilde T=1.116\cdot10^{22}$\,s$^{-1}$.

\section{The temperature of the bit: scale, state, bath}
The energy $E_0$ converts directly into a \emph{temperature}. From $E_0=k_B T_0$,
\[
T_0=\frac{E_0}{k_B}=\frac{1.177\cdot10^{-12}\,\text{J}}{1.381\cdot10^{-23}\,\text{J/K}}
   =8.53\cdot10^{10}\,\text{K}
\]
(check via $1$\,eV$=11\,604.5$\,K: $7.348\cdot10^6\,\text{eV}\cdot11\,604.5=8.53\cdot10^{10}$\,K), about
$85$ billion kelvin. $T_0$ is the temperature at which the thermal energy $k_B T_0$ reaches the
\emph{structure energy} $E_0$ of the bit.

\paragraph{Does a bit even have a temperature?} Strictly: no, not \emph{one}. A single bit in a definite
state ($0$ or $1$) has zero entropy and no distribution -- temperature is an ensemble/reservoir quantity,
not assigned to a single fixed degree of freedom. What one \emph{can} sensibly assign to a bit are
\emph{three distinct} temperature-like quantities that must not be confused. First, the \emph{state} or
occupation temperature: if the bit has two levels split by $E_0$, the occupation ratio
$p_1/p_0=e^{-E_0/k_B T}$ defines an effective temperature. It runs \emph{against} intuition: maximally
mixed ($p_0=p_1=\tfrac12$, a full one bit of entropy) means $T_\text{eff}\to\infty$, definite ($p\to0$ or
$1$, zero entropy) means $T_\text{eff}\to0$, population inversion even $T_\text{eff}<0$. The
\emph{most informative} bit is thus the infinite-temperature state; cooling drives it to a definite state
and \emph{destroys} the information. (A purely logical, degenerate bit -- both values same energy -- has
no occupation temperature at all.) Second, the \emph{scale} $T_0=E_0/k_B$: not \enquote{where the bit
sits}, but the characteristic melting/activation point from the energy gap, the natural temperature
\emph{unit} of the bit. Third, the \emph{bath} temperature that enters the erase cost -- and that belongs
to the \emph{environment}, not the bit.

The scale $T_0$ marks the thermal robustness threshold: above it ($k_B T>E_0$) the thermal energy
overwhelms the structure, and the heat washes out the distinction -- the winding or mode that the bit
\emph{is}: the bit \enquote{melts}. Below it ($k_B T<E_0$) it is robust. This is a statement about the
\emph{scale}, not about a state.

\paragraph{Which temperature when erasing?} The Landauer cost $Q=k_B T\ln2$ depends on the temperature of
the \emph{bath} into which the erased entropy flows as heat -- on the choice of environment, not on the
bit. There is no universal \enquote{correct} value. For a real computer it is the operating temperature,
in practice \emph{room temperature} ($\sim300$\,K): $Q=k_B\cdot300\cdot\ln2=18$\,meV. That is the honest
answer for real computing. UIFT instead uses $T_\text{CMB}=2.7$\,K -- not absolute zero, but the
\emph{coldest natural heat sink}; this gives the minimum possible erase heat in the present universe
($Q=0.16$\,meV, nothing colder without active cooling). Legitimate as a cosmological floor, but still just
\emph{one} bath choice, not a bit property. $T_0$ finally is \emph{not} a bath: inserting it gives the
limiting case \enquote{bath as hot as the bit itself}, $Q=E_0\ln2=0.69\,E_0$ -- a reference, not an
erasure.

\enquote{The} temperature of a bit therefore depends on its \emph{scale}: a generic bit of length $L$
carries energy $E=\hbar c/L$ and hence scale $T=\hbar c/(k_B L)$. Each bit scale has its own;
$8.5\cdot10^{10}$\,K is the one belonging to the characteristic $E_0$. A larger $L$ -- a \enquote{softer}
bit -- has a lower threshold.

So $T_0$ stands alongside mass and time as a third unit clothing of the \emph{same} energy. With one
bookkeeping constant each:
\[
m_0 c^2 \;=\; E_0 \;=\; k_B T_0 \;=\; \hbar/\tilde T,
\]
i.e.\ $/c^2\to$ mass $m_0=1.310\cdot10^{-29}$\,kg, $/\hbar\to$ time $\tilde T=8.958\cdot10^{-23}$\,s (the
dual time from $\tilde T\cdot m=1$), $/k_B\to$ temperature $T_0=8.53\cdot10^{10}$\,K. Beneath all four
sits the dimensionless count $\log_2 M$. One energy, four clothings.

The temperature is not merely \enquote{also there} but tied directly to the dynamic dual time. The
\emph{thermal time} at $T_0$ is $\hbar/(k_B T_0)$; inserting $T_0=E_0/k_B$ gives $\hbar/E_0=\tilde T$ --
the thermal time at $T_0$ \emph{is} the dual time from $\tilde T\cdot m=1$ (numerically identical to
$10^{-38}$\,s, \texttt{bit\_temperatur.py}). This is the KMS / \enquote{thermal time} correspondence
(Connes--Rovelli): temperature is an inverse (imaginary) time, and it lands without extra assumption
exactly on the FFGFT dual time. On the bookkeeping axis $k_B$ is for temperature exactly what $\hbar$ is
for time.

Finally, $T_0$ joins the two \enquote{energies} of a bit -- what it \emph{is} ($E_0$) and what its
\emph{erasure} costs ($k_B T\ln2$) -- into \emph{one} scale. The Landauer erase heat relative to the
self-energy is simply the temperature ratio,
\[
\frac{Q}{E_0}=\frac{k_B T\ln2}{k_B T_0}=\frac{T}{T_0}\,\ln 2.
\]
At room temperature: $\tfrac{300}{8.53\cdot10^{10}}\ln2=2.4\cdot10^{-9}$, hence
$E_0/Q=(T_0/T)/\ln2\approx4\cdot10^{8}$ -- room temperature sits about $8.5$ decades below $T_0$, and
\emph{that} is why the static self-energy ($\sim7.3$\,MeV) and the erase heat ($\sim18$\,meV) lie a good
eight orders of magnitude apart. At $T=T_0$, by contrast, erasure costs $E_0\ln2=0.69\,E_0$: at its own
temperature point, \emph{being} and \emph{erasing} are the same order of magnitude. This ratio $T/T_0$ is
the \emph{only} setup-independent statement: a bit has no absolute temperature, only its bath temperature
relative to its own scale. The static ($E_0,T_0$) and thermodynamic ($k_B T\ln2$) sides are thus joined by
the single ratio -- no jump, one scale.

Placement: $T_0\approx8.5\cdot10^{10}$\,K is the MeV/nuclear scale -- the temperature order of nuclear and
hadronic physics: stellar cores (solar core $\sim1.5\cdot10^7$\,K) and the quark-gluon plasma in heavy-ion
collisions (QCD transition $\sim1.7\cdot10^{12}$\,K at $150$\,MeV). The \enquote{hot early universe}
reading is deliberately left out -- it is a $\Lambda$CDM assumption that FFGFT does not share (FFGFT's CMB
structure follows from the T$^4$ geometry, not from a hot early phase). Honestly:
$T_0=E_0/k_B$ is the \emph{standard} energy-temperature assignment; what is FFGFT-specific is only that it
coincides with the dual time via $\tilde T\cdot m=1$ -- a clean consistency, not a new derivation. And
because the absolute erase temperature is the bath's, a $\xi$ tied to a particular bath ($T_\text{CMB}$ or
room temperature) would depend on that choice -- exactly why the CMB-bound $\xi_\text{UIFT}$ is a
contingent, different quantity while FFGFT's geometric $\xi$ stays bath-\emph{independent} (\S9).

\paragraph{$E_0$ is a rung, not the floor.} The scale $E_0$ is not the bottom end of the framework but a
\emph{rung} of a fractal ladder. The geometric floor lies at the sub-Planck length
$\xi\,\ell_P\approx2.2\cdot10^{-39}$\,m (Doc.~163) -- the smallest length below which spacetime carries no
well-defined degrees of freedom. Because $\xi\ell_P$ is a factor $\xi$ \emph{below} the Planck length,
this floor stands, in the \enquote{hard} clothings, a factor $1/\xi=7500$ \emph{above} Planck:
\begin{center}
\begin{tabular}{@{}lll@{}}
Length & $\xi\ell_P$ & $2.16\cdot10^{-39}$\,m\\
Time/duration & $\xi t_P$ & $7.19\cdot10^{-48}$\,s\\
Frequency & $1/(\xi t_P)$ & $1.39\cdot10^{47}$\,Hz\\
Energy & $E_P/\xi$ & $9.16\cdot10^{22}$\,GeV\\
Temperature & $T_P/\xi$ & $1.06\cdot10^{36}$\,K\\
Mass & $m_P/\xi$ & $1.63\cdot10^{-4}$\,kg\\
\end{tabular}
\end{center}
This floor is the \emph{smallest-length} edge, hence in energy/temperature/frequency the \emph{maximum}
(UV): \enquote{floor} means finest resolution, not smallest energy. The bit rung $E_0$ ($7.348$\,MeV,
$26.85$\,fm, $T_0=8.5\cdot10^{10}$\,K) sits about $25$ decades \emph{coarser}. Between them the fractal
runs inward, factor $1/\xi=7500$ per main level (Doc.~146), with $\approx13$ $\sqrt2$ sub-steps each -- and
$\log_2 7500\approx12.87$ is \emph{exactly} the information ceiling $I_1$ per degree of freedom from \S5.
The number of fractal sub-steps per level \emph{is} the maximum bit count; the capacity of the bit and the
fineness of the fractal are the same number.

\paragraph{Landauer.} Erasing one bit is logically irreversible: two states are mapped to one, the entropy
drops by $k_B\ln2$. By the second law this entropy must flow to the environment as heat, $Q\ge k_B T\ln2$
(Landauer 1961). Here $T$ is the temperature of the \emph{bath} into which the erasure goes -- not the
bit's.

\paragraph{The problem.} And here it jams. The framework now supplies a whole ladder of \emph{definite}
temperatures -- the rung $T_0=8.5\cdot10^{10}$\,K, the floor $T_{\max}=1.1\cdot10^{36}$\,K, each fractal
level its own in between, over $25$ decades. \emph{None} of them is \enquote{the} Landauer temperature.
Insert an intrinsic scale into $Q=k_B T\ln2$ and you get no real erasure but a limiting case or nonsense:
with $T_0$ the reference $Q=0.69\,E_0$, with $T_{\max}$ an absurdly large heat, with $T_\text{CMB}$ the
UIFT number (\S9). The real erasure takes the temperature of the \emph{actual environment} (computer:
$\sim300$\,K). That is the problem: despite a ladder of definite intrinsic temperatures, the erase
temperature remains the bath's, and the only invariant is the \emph{ratio}
$Q/E=(T/T_\text{rung})\ln2$ -- bath relative to the respective rung. A bit has a definite energy and, per
rung, a definite scale temperature; a definite \emph{erase price} it acquires only with an environment.

\section{Consequences}
\emph{The faithful description is quantum, not bit-based.} Because the Hilbert space is a bijection
(Doc.~230) but the classical bit count a projection: FFGFT's lossless description is the
amplitude-plus-phase representation (quantum information); the Shannon bit count is its lossy shadow. An
information primitive needs magnitude \emph{and} phase.

\emph{The open datum is relational information.} The one irreducible open quantity (Doc.~282) is the
holonomy phase $\theta$ -- path-dependent, relational, not a local bit content. The residue the bit
projection discards \emph{is} the phase; and the phase \emph{is} where the open physics lives. The
fundamental information sits in the connection: the electromagnetic field $A_\mu$ is that connection, the
photon its quantum -- a conceptual example, not one numerically identical to $\theta$. \enquote{It from
bit} becomes \enquote{it from connection}.

\section{Honest boundaries}
\textbf{Structural diagnosis, not a new derivation} (P35). Doc.~255/257 already had Landauer/Vopson; new
is the magnitude/phase decomposition and its consequences. The \emph{one exact} point:
$E_\text{bit}=\hbar c/L$ is a photon energy.

\textbf{Photon is a gauge boson, neutrino a fermion.} The photon analogy holds at the level of
masslessness/dispersion, not as an identity (spin 1 vs.\ $\tfrac12$); the link to the connection belongs to the gauge field $A_\mu$ (the photon being its quantum), only
in the photon sector. The photon's phase (polarization/gauge holonomy) is the same \emph{category} as $\theta$,
not the same number. The photon sits outside the lepton mass operator ($m=0$ reference).

\textbf{$\log_2(1/\xi)\approx12.87$ is not a magic number.} The content is finiteness and $\xi$
anchoring, not the value; the exact form ($1/\xi$ vs.\ $2/\xi$, direction/flavour factor) depends on the
mode counting. Direction: geometry $\to$ information, not the reverse.

\textbf{Energy side: $\xi$-anchored, the value match is empirical.} Theoretically
$\xi$-/geometry-fixed ($E_0=7.398$~MeV, Doc.~011); for the SI numbers the empirical value
$\sqrt{m_e m_\mu}=7.348$~MeV is used so the comparison with measured quantities fits -- from it
$m_0,T_0,\tilde T,1/\tilde T$ follow exactly ($\hbar,c,k_B$ are mere bookkeeping, no deviation). The only
difference is the dimensionless theory-vs-empirics match $\alpha/\xi$ vs.\ $m_e m_\mu$ ($+0.68\%$), not a
conversion.

\textbf{Bit thermodynamics: computed, over the bit's own scale $T_0$.} The bit has the full Landauer
computation -- erase cost $Q=k_B T\ln2$, relative to the self-energy $Q/E_0=(T/T_0)\ln2$ (\S7); at
$T=T_0$, $Q=E_0\ln2=0.69\,E_0$, the intrinsic quantity is the bit entropy $\ln2$. No CMB needed. The UIFT
parameter $\xi_\text{UIFT}=k_B T_\text{CMB}\ln2/(m_e c^2)\approx3\cdot10^{-10}$ (Doc.~013/243) is the same
Landauer form with foreign cosmological scales ($T_\text{CMB}$ instead of $T_0$, $m_e$ instead of $E_0$)
-- a different quantity, not FFGFT's geometric $\xi=1.3\cdot10^{-4}$. A different scale choice, not a
failure.

\textbf{Neutrino argument: speculative.} It still rests on the speculative degeneracy (Doc.~007) -- but it
no longer carries the thesis; the photon does.

\section{Balance: not whether, but which channel}
The magnitude/phase work shifts the information question from \enquote{is information fundamental?} to
\enquote{which channel?}. Vopson/Landauer are the magnitude channel -- real, derived from the $T^4$
geometry, but blind to the phase. The phase channel (mixing, holonomy, $\theta$) is independent,
relational information; in signal processing it is the one that carries the structure. The cleanest
anchor is the photon: exactly massless, FFGFT's information unit $E_\text{bit}=\hbar c/L$ itself, and the quantum of
the $U(1)$ connection -- conceptually \enquote{it from connection}. The \emph{amount} of information is
absolute and anchored in $\xi$ ($I_1=\log_2(1/\xi)\approx12.87$ bit); energy and capacity are its
dimensional dressings. Only the geometry holds both channels -- honestly bounded as diagnosis, not a new
derivation.

